\section{Lecture 9 (September 29th)}
\begin{defi}
(Spin-$\dfrac{1}{2}$ Dirac Fields) The Dirac action was given as
\[S_{D}=\int d^{4}x\,\bar{\psi }(i\cancel{\partial }-m)\psi \]
The canonical quantisation or ETCR were give as
\[\hat{\psi }(x),\hat{\pi }(x)=i\hat{\psi }^{\dagger}(x)\longrightarrow 
\begin{cases}
\{\hat{\psi }_{\alpha }(t,\vec{x}),\hat{\psi }_{\beta }(t,\vec{y})\}=\{\hat{\pi }_{\alpha }(t,\vec{x}),\hat{\pi }_{\beta }(t,\vec{y})\}=0\\
\{\hat{\psi }_{\alpha }(t,\vec{x}),\hat{\pi }_{\beta }(t,\vec{y})\}=i\delta _{\alpha \beta }^{(4)}(\vec{x}-\vec{y})
\end{cases}\]
where the symmetriser comes from the spin-statistics theorem (anti-commutation for fermions). The equations of motion are promoted from the Hamiltonian as
\[\hat{H}=\int d^3\vec{x}\,\hat{\psi }^{\dagger}(-i\gamma^{0}\vec{\gamma }\cdot \vec{\nabla }+\gamma ^{0}m)\hat{\psi }\]
The Hamilton's equation become 
\[\dfrac{d \hat{\psi }}{d t}=-i[\hat{\psi },\hat{H}]\longrightarrow (-i\cancel{\partial }-m)\hat{\psi }=0\]
The solutions to the quantised Dirac equation look like
\[\hat{\psi } (x)=\int \dfrac{d^3\vec{p}}{\sqrt{(2\pi )^32E_{p}}}\sum ^{2}_{i=1}\Big(\hat{b}^{i}(\vec{p})u^{i}_{+}(\vec{p})e^{ip\cdot x}+\hat{d}^{i}_{\vec{p}}^{*}u_{-}^{i}(\vec{p})e^{-ip\cdot x}\Big)\]
while the conjugate field looks like 
\[\hat{\pi }(x)=i\int \dfrac{d^3\vec{p}}{\sqrt{(2\pi )^3}2E_{p}}\sum ^{2}_{i=1}\Big(\hat{b}^{i}(\vec{p})^{\dagger}u_{+}^{i}(\vec{p})^{\dagger}e^{-ip\cdot x}+\hat{d}^{i}(\vec{p})u^{i}_{-}(\vec{p})^{\dagger}e^{ip\cdot x}\Big)\]
Recall that after we found the fields, we wanted the mode operators in terms of the fields and we obtain
\[\hat{b}^{i}(\vec{p})=\int \dfrac{d^3\vec{x}}{\sqrt{(2\pi )^32E_{p}}}\,u^{i}_{+}(\vec{p})^{\dagger}\hat{\psi }(x)e^{-ip\cdot x}\]
and
\[\hat{d}^{i}(\vec{p})^{\dagger}=\int \dfrac{d^3\vec{x}}{\sqrt{(2\pi )^32E_{p}}}\,u_{-}^{i}(\vec{p})^{\dagger}\hat{\psi }(x)e^{ip\cdot x}\]
(this can be proved through using $u^{i}_{\pm}(\vec{p})^{\dagger}u^{j}_{\pm}(\vec{p})=2E_{p}\delta ^{ij}$ and $u_{\pm}^{i}(\vec{p})u^{j}_{\mp}(-\vec{p})=0$). The anti-commuation relations for $\hat{b}^{i}$ and $\hat{d}^{i}$ should be
\[\{\hat{b}^{i}(\vec{p}),\hat{b}^{j}(\vec{q})^{\dagger}\}=\{\hat{d}^{i}(\vec{p}),\hat{d}^{j}(\vec{q})^{\dagger}\}=\delta^{ij}\delta ^{(3)}(\vec{p}-\vec{q})\]
which is exactly the infinitely many sets of creation and annihilation operators for a fermionic harmonic oscillator.
\end{defi}
\vspace{2ex}
\begin{defi}
(Construction of the Hilbert/Fock space) The vaccum state is defined using the annihilator operator 
\[\hat{b}^{i}(\vec{p})|0\rangle =\hat{d}^{i'}(\vec{p})|0\rangle =0\]
The first excited states look something like
\[\hat{b}^{i}(\vec{p})^{\dagger}|0\rangle =|(\vec{p},i)\rangle \qquad\&\qquad \hat{d}^{i'}(\vec{p})^{\dagger}|0\rangle =|(\vec{p}',i')\rangle \]
The multiple particle states trivially follow as
\[|(\vec{p},i_1),\ldots ,(\vec{p}_{M},i_{N}),(\vec{p}_{1}',i'_{1}),\ldots ,(\vec{p}'_{N},i_{N}')\rangle =\hat{b}^{i_1}(\vec{p}_{1})^{\dagger}\ldots \hat{d}^{i_{N}'}(\vec{p}_{N}')^{\dagger}|0\rangle \]
\end{defi}
\vspace{2ex}
\begin{thm}
(Pauli exclusion principle) This naturally follows from
\[\Big(\hat{b}^{i}(\vec{p})^{\dagger}\Big)^2=\Big(\hat{d}^{i'}(\vec{p}\,')\Big)^2=0\]
and the general multiple particle state vanishes when exists $(m,n)$ such that
\[(\vec{p}_{m},i_{m})=(\vec{p}_{n},i_{n})\qquad\mathrm{or}\qquad (\vec{p}_{m}\,',i_{m}')=(\vec{p}_{n}\,',i_{n}')\]
\end{thm}
\vspace{2ex}
\begin{defi}
The momentum and Hamiltonian in terms of the mode operators are
\[\hat{P}^{\mu }=\int d^3\vec{x}\,\hat{T}^{0\mu }=-\int d^3\vec{x}\,i\hat{\bar{\psi }}\gamma ^{0}\partial ^{\mu }\hat{\psi }\]
and
\[\hat{H}=i\int d^3\vec{x}\,\hat{\psi }^{\dagger}\partial _{t}\hat{\psi }\]
resulting in
\[\begin{align*}
\hat{H}=&i\int \dfrac{d^3\vec{p}}{2E_{p}}(-iE_{p})\sum _{i,j=1}^{2}2E_{p}\delta ^{ij}\Big(\hat{b}^{i}(\vec{p})^{\dagger}\hat{b}^{j}(\vec{p})-\hat{d}^{i}(\vec{p})\hat{d}^{j}(\vec{p})^{\dagger}\Big)\\
=&\int d^3\vec{p}\,E_{p}\sum ^{2}_{i=1}\Big(\hat{b}^{i}(\vec{p})^{\dagger}\hat{b}^{i}(\vec{p})-\hat{d}^{i}(\vec{p})\hat{d}^{i}(\vec{p})^{\dagger}\Big)\\
=&\int d^3\vec{p}\,E_{p}\sum ^{2}_{i=1}\Big(\hat{b}^{i}(\vec{p})^{\dagger}\hat{b}^{i}(\vec{p})+\hat{d}^{i}(\vec{p})^{\dagger}\hat{d}^{i}(\vec{p})\Big)-\int d^3\vec{p}\,E_{p}\delta ^{(3)}({\bf 0})
\end{align*}\]
The momentum operator is naturally 
\[\begin{align*}
\hat{\vec{P}}=&-i\int d^3\vec{x}\,\hat{\psi }^{\dagger}\nabla \hat{\psi }\\
=&\int d^3\vec{p}\,\vec{p}\sum ^{2}_{i=1}\Big(\hat{b}^{i}(\vec{p})^{\dagger}\hat{b}^{i}(\vec{p})+\hat{d}^{i}(\vec{p})^{\dagger}\hat{d}^{i}(\vec{p})\Big)
\end{align*}\]
There are also other important conserved charges! The $U(1)$ symmetry becomes the particle numbers ($+1$ for $\hat{b}^{\dagger}$ and $-1$ for $\hat{d}^{\dagger}$). The Noether current is $j^{\mu }=\bar{\psi }\gamma ^{\mu }\psi $ leading to the conserved quantity
\[\hat{N}=\int d^3\vec{x}\,\hat{j}^{0}=\int d^3\vec{x}\,\hat{\psi }^{\dagger}\hat{\psi }\]
and
\[\begin{align*}
\hat{N}=&\int d^3\vec{p}\,\sum ^{2}_{i=1}\Big(\hat{b}^{i}(\vec{p})^{\dagger}\hat{b}^{i}(\vec{p})+\hat{d}^{i}(\vec{p})\hat{d}^{i}(\vec{p})^{\dagger}\Big)\\
=&\int d^3\vec{p}\,\sum ^{2}_{i=1}\Big(\hat{b}^{i}(\vec{p})^{\dagger}\hat{b}^{i}(\vec{p})-\hat{d}^{i}(\vec{p})^{\dagger}\hat{d}^{i}(\vec{p})\Big)+N_0
\end{align*}\]
Also, the normal Lorentz symmetries give the generalised angular momenta. For the Dirac field, 
\[\mathcal{J}^{\mu \nu \rho }=\underbrace{x^{\rho }T^{\mu }-x^{\nu }T^{\mu \rho }}_{\mathrm{orbital}}+\underbrace{\bar{\psi }\gamma ^{\mu }\Big(\dfrac{i}{2}\gamma ^{\nu \rho }\Big)\psi}_{\mathrm{spinorial}=\mathcal{S}^{\mu \nu \rho }} \]
We are compelled to think that
\[\begin{cases}
\hat{J}^{\mu \nu }=\int d^3\vec{x}\,\hat{\mathcal{J}}^{0\mu \nu }\\
\hat{S}^{\mu \nu }=\int d^3\vec{x}\,\hat{\mathcal{S}}^{0\mu \nu }
\end{cases}\]
is the second part generalised spin? There are a lot of subtleties, addressed in the homework problems.
\end{defi}
\vspace{2ex}
\begin{rmk}
The vaccum energy $E$ can be seen as the sum of all ground state energies of the fermionic oscillators. They are therefore negative and divergent. Let us interpret the ground state recalling that 
\[\{\hat{d},\hat{d}^{\dagger}\}\qquad\&\qquad \hat{d}|0\rangle =0\]
where the first is symmetric.

\bigskip

The first perspective is that $\hat{d}^{\dagger}$ and $\hat{d}$ create and annihilate positive energy states. Therefore, $|0\rangle $ can be seen as the vaccum without any excited positive energy states. Also, $\hat{b}^{\dagger}$ and $\hat{d}^{\dagger}$ would be creating particles and anti-particle states (positive). 

\bigskip

The second perspective is that $\hat{d}$ and $\hat{d}^{\dagger}$ create and annihilate negative energy states. Then, the state $|0\rangle $ would be interpreted as a state with fully filled negative energy states, exactly the "Dirac sea". Also, $\hat{d}^{\dagger}$ would create a hole in the Dirac sea. 
\end{rmk}
\vspace{2ex}
\begin{defi}
(Normal ordering) Again, the normal ordering operator sends all annihilation operators to the right, giving
\[\begin{align*}
:\hat{H}:=&\int d^3\vec{p}\,E_{p}\sum ^{2}_{i=1}\Big(\hat{b}^{i}(\vec{p})^{\dagger}\hat{b}^{i}(\vec{p})+\hat{d}^{i}(\vec{p})^{\dagger}\hat{d}^{i}(\vec{p})\Big)\\
:\hat{\vec{P}}:=&\int d^3\vec{p}\,\vec{p}\sum ^{2}_{i=1}(\ldots )\\
:\hat{N}:=&\int d^3\vec{p}\,\sum ^{2}_{i=1}\Big(\hat{b}^{i}(\vec{p})^{\dagger}\hat{b}^{i}(\vec{p})-\hat{d}^{i}(\vec{p})^{\dagger}\hat{d}^{i}(\vec{p})\Big)
\end{align*}\]
Note that $:\hat{a}\hat{a}^{\dagger}:=\hat{a}^{\dagger}\hat{a}$ for bosons whereas $:\hat{b}\hat{b}^{\dagger}:=-\hat{b}^{\dgger}\hat{b}$ for fermions.
\end{defi}
\vspace{2ex}

